\section{{Introduction}}

Much of what people ask about public services has an answer that holds only under a condition. Can I bring my other child on this visa? Yes, if that child also has or is applying for a Child Student visa. Do I have to keep payroll records? Yes, for at least six years, if they were created after a certain date. An answer stripped of its condition is not merely incomplete; for the person whose situation does not meet the condition, it is wrong, and it is stated with the confidence of a correct one.

Retrieval-augmented generation (RAG) \cite{lewis2020rag} is now a common way to answer such questions from official guidance. A RAG pipeline never reads a guidance page whole. It reads \emph{chunks}: pieces of fixed size cut from each page, retrieved by similarity to the question, often filtered or reranked, and then passed to a language model. Each of these steps can separate an answer from the condition that limits it, and each can do so silently, since the answer that remains is fluent and partly right.

At the same time, a large body of work makes RAG pipelines more selective about evidence, because retrieved passages can be irrelevant or false \cite{xu2024survey}. Relevance graders drop passages judged unhelpful \cite{yan2024crag}; conflict resolvers detect contradicting passages and keep the more credible one \cite{conflictrag2026}; consolidation methods keep what they judge reliable \cite{wang2025astute}. These modules are designed and evaluated where removing evidence is the right thing to do, such as misinformation. They are rarely evaluated where the passage that looks least relevant, or most contradictory, is an exception to the rule beside it.

This paper asks a measurement question: \emph{where, in a chunked RAG pipeline, are the conditions of an answer lost?} We answer it on ConditionalQA \cite{sun2022conditionalqa}, whose answers are annotated by people with the conditions under which they hold, turned into an open-retrieval task over all \nPages{} of its gov.uk pages. Because the human annotation identifies every condition as an exact element of its page, the first three stages of the pipeline can be measured directly, with no language model in the loop: whether chunking separates a condition from its answer, whether retrieval finds it, and whether an evidence-reducing module keeps it. The fourth, generation, is measured with the benchmark's official metric, under a single generator prompt shared by every system so that systems differ only in the evidence they pass on.

Our main findings are as follows.
\begin{enumerate}[(1)]
\item \textbf{Chunking separates conditions from answers.} With 150-word chunks, the gold evidence of \sepCond\% of conditional questions spans more than one chunk, against \sepUncond\% of unconditional ones; for \condOnly\% of conditional questions some condition sits in a chunk that holds nothing else the answer needs.
\item \textbf{Retrieval is not biased against conditions, except those chunked apart.} Condition chunks are retrieved as often as answer chunks (\recCondEight\% and \recAnsEight\% at $k=8$), but chunks that hold only a condition are retrieved less (\recCondOnlyEight\%).
\item \textbf{Evidence-reducing operators are blind to conditions.} A reranker, a contradiction filter and a ConflictRAG-style resolver remove retrieved condition chunks at the same rate as answer chunks; the contradiction filter fires on \nliChanged\% of questions about a single, consistent body of guidance, because an exception reads as a contradiction of its rule. On RAMDocs \cite{wang2025madam}, where discarding is the point, the same operators do not single out misinformation either.
\item \textbf{Generation is where most conditions are lost.} With every condition in context, the two local generators state \recallGoldQwen\% and \recallGoldLlama\% of the gold conditions as conditions on their answer, and the largest of three larger hosted models \hRecallNemo\%; the official metric, meanwhile, rewards a generator for rarely stating conditions at all.
\end{enumerate}

End to end, two local generators and three larger hosted ones agree on the order of these losses (Sections~\ref{sec:endtoend} and~\ref{sec:hosted}). Losing an isolated condition chunk from otherwise complete evidence costs every generator measurably; retrieval at an equal budget costs about a point; the operators, despite their deletions, change little for small generators, which have already lost most conditions in generation; and the larger generators, which state more conditions, lose more at every stage before generation. We also test a remedy aimed at the retrieval side, \emph{qualifier-aware context completion}, which adds the same-page chunks that a cross-fitted scorer predicts state a condition. It helps none of the five generators, and its scorer neither beats an untrained reranker nor transfers to a separately annotated corpus. We report it as a negative result: the lever is the generator, not the context.

The contributions of this paper are: (i) a stage-wise, annotation-grounded measurement of condition loss in chunked RAG; (ii) the first evaluation of evidence-reducing and conflict-handling operators on ConditionalQA, contrasted with the same operators on RAMDocs; (iii) a pre-registered test of qualifier-aware context completion, a cross-fitted remedy trained only on the benchmark's existing annotation, which does not help; and (iv) all code, the fixed question splits, every model output and an analysis plan committed before the first evaluation run.
